\section{The preservation property}
\label{sec:preservation}
The protocol has a single composition criterion: every allowed change must
preserve the bounds and identity of existing commitments. Three invariants
make this criterion precise.

\subsection{Capacity is allocated once}
Let $B_D$ be a delegation root's ceiling, $b_v$ the ceiling of a leaf and $U_v$
the unused remainder of an internal slot. Subdivision preserves
\begin{equation}
 \sum_{\mathrm{leaves}\ v}b_v+\sum_{\mathrm{internal}\ v}U_v=B_D,
 \qquad b_v,U_v\geq0. \label{eq:delegation}
\end{equation}
It moves capacity from a parent remainder into a child, without changing the
sum. Internal slots cannot execute; selection and sealing change no ceiling.
Declared effect and reader inclusion, decreasing expiry and decreasing depth
extend transitively along the allocation tree. With the bound one-invocation
capabilities and native charge enforcement, admitted leaf charges fit those
ceilings.

For a swarm lineage, let $A_i$ be the set of uniquely identified allocations
in $G_i$, with fixed ceiling $b(a)$ and original declared pool $B_G$.
Every checked extension satisfies
\begin{equation}
 A_i\subseteq A_{i+1},\quad
 b_{i+1}(a)=b_i(a)\ (a\in A_i),\quad
 \sum_{a\in A_{i+1}}b(a)\leq B_G. \label{eq:growth}
\end{equation}
A single protected lineage per declared pool, with a serialized head,
establishes the domain of this bound. Because earlier allocations are
retained, $\bigcup_{i=0}^n A_i=A_n$: the latest bound also covers every allocation
ever installed. These are capacity bounds; deposited assets have a separate
accounting invariant.

\subsection{Funds follow their own allocation}
For each accounted endowment $u$, fix deposited funds $D_u$ with no further
deposits into that accounting cohort. Let $A_u$ be its available balance,
$L_u$ its locked funds, and $P_u,R_u$ its paid and refunded outflows. Then
\begin{equation}
 D_u=A_u+L_u+P_u+R_u,\qquad A_u,L_u,P_u,R_u\geq0.
 \label{eq:conservation}
\end{equation}
Reservation moves funds from $A_u$ to $L_u$; payment or refund moves an
allocation's locked funds once to the corresponding outflow. Each transition
preserves the sum and nonnegativity. Returned cash remains an outflow in this
cohort. A later deposit is accounted as another endowment; it neither resets
these accounts nor releases their still-locked allocations.

An earned allocation remains in $\{\textsf{Payable},\textsf{Paid}\}$.
A transition on another allocation leaves it unchanged. Thus a parent refund
cannot consume an earned child's reserve. Eventual withdrawal also requires
the beneficiary's key, available evidence, transaction inclusion and a
functioning asset.

\subsection{Compose through the original invocation}
\begin{proposition}[Preservation under evolving work]
Assume the protected allocation, graph, receiver-custody and settlement domains
of Section~\ref{sec:sovereignty}, the authenticated bindings of
Section~\ref{sec:contract}, and the transitions of Section~\ref{sec:execution}.
Begin with validly provisioned roots, verified initial graphs and funded
accounts. Within those declared domains, every finite interleaving of valid
subdivision, sealing, graph growth, native execution, recovery and settlement
preserves the following properties:
\begin{enumerate}
 \item Declared delegation and swarm capacity remain within their original
       bounds; paid plus refunded funds never exceed deposited backing.
 \item A sealed invocation retains its selected receiver, capability, input,
       effect and request. Graph growth and replay grant it no additional
       native invocation right.
 \item A recorded earned child claim retains its original terms and remains
       payable or paid through any parent failure, completion or refund.
\end{enumerate}
\end{proposition}

\begin{proof}
Initially, a new root has all its capacity in one leaf, a verified graph fits
its declared pool, and a fresh funded endowment has
$(A_u,L_u,P_u,R_u)=(D_u,0,0,0)$. No execution right has been consumed or payment
earned. The properties hold in this initialized state.

Induct over the interleaving. A subdivision changes only the parent's remainder
and its new child, preserving (\ref{eq:delegation}). A selection changes no
capacity. Sealing freezes the selected binding before evidence escapes.

A graph extension preserves every earlier allocation and checks
(\ref{eq:growth}) before committing its new head. The old continuation keeps its
resource identity. Exact historical lookup preserves the artifacts used by
an old request; a new graph signature cannot cause native custody to acquire
that resource as a second invocation.

Admission compares the request with the sealed permit, native capability,
graph task, route and treaty. The funded agreement authenticates that same
request. A substituted object either fails the binding or requires a different
authorized commitment. Native recovery resolves the original operation, hold
and continuation claim. Refusal, expiry and revocation can prevent fresh work;
they do not create a replacement invocation from the old commitment.

Financial transitions preserve (\ref{eq:conservation}). Their allocation
identity confines each transition to its own reserved funds. A parent
transition cannot change the child; a child transition after acceptance stays
inside $\{\textsf{Payable},\textsf{Paid}\}$.

These cases cover the allowed transitions. They remain valid when interleaved
because additive growth retains the resources on which old execution depends,
and each scarce resource is serialized by its own custodian.
\end{proof}

The result explains why the construction needs no atomic transaction spanning
all owners. An incomplete handoff can strand a commitment, but cannot use
absence of a later record as permission to replace it. Safety is obtained by
retaining the earlier fact and resolving it under its original authority.

This is a conditional protocol argument. The accompanying Lean development
checks abstract financial conservation and earned-state lemmas; native tests
exercise concrete admission and recovery boundaries. Their relationship to the
Rust and Solidity implementations is tested rather than established by a
machine-checked refinement proof.

Exclusive allocation is a necessary premise. Two receivers can each see a valid
promise against one deposit while neither sees the other promise. Signatures
alone cannot distinguish that execution from two individually affordable
exchanges. A protected allocator, trusted ledger or settlement chain must
supply the shared allocation fact. Chio makes that custodian explicit at each
handoff.
